\documentclass[10pt,a4paper]{amsart}
\usepackage[margin=19mm]{geometry}
\usepackage{amsmath,amssymb,mathtools}
\usepackage[scr=rsfs,cal=euler]{mathalfa}
\usepackage{xfrac}
\usepackage[unicode,hidelinks,bookmarks=false]{hyperref}
\newcommand{\ie}{i.e.\ }
\newcommand{\cf}{cf.\ }
\newcommand{\resp}{respectively\ }
\newcommand{\Subsection}{\subsection}
\providecommand{\nobreakdash}{\nobreak\hbox{-}}
\setlength{\emergencystretch}{2em}
\multlinegap0pt
\usepackage{bm}
\usepackage{bbm}
\usepackage{dsfont}
\usepackage{xspace}
\usepackage{cancel}
\usepackage{mathtools}
\usepackage{amsthm}
\makeatletter
\@ifundefined{thm}{\newtheorem{thm}{Thm}}{}
\@ifundefined{lem}{\newtheorem{lem}[thm]{Lemma}}{}
\makeatother
\title[Spectral analysis of $q$-deformed unitary ensembles with the]{Spectral analysis of $q$-deformed unitary ensembles with the Al-Salam--Carlitz weight}
\author{Sung-Soo Byun, Yeong-Gwang Jung, Jaeseong Oh}
\date{Source: arXiv:2507.18042v1 (CC BY 4.0)}
\begin{document}
\maketitle

\noindent\emph{Source and licence.} Spectral analysis of $q$-deformed unitary ensembles with the Al-Salam--Carlitz weight. Version arXiv:2507.18042v1 (CC BY 4.0), distributed under the Creative Commons Attribution 4.0 International licence, as recorded in the arXiv licence metadata for this exact version and shown on the abstract page https://arxiv.org/abs/2507.18042v1. Full rights notice: licenses/arxiv-2507.18042-CC-BY-4.0.txt.\par\smallskip

\noindent\textit{Editorial scope.} Counted unit: the Lem whose source label is \texttt{lem:H asymptotic} in the article (source0.tex lines 1520--1537, source label \texttt{lem:H asymptotic}), delivered together with the article's own complete original proof of it (source0.tex lines 1538--1583, one \texttt{proof} environment of 46 lines). No mathematics is written, repaired or re-derived: statement and proof are the article's own text line for line. Required context retained uncounted: def of H (lines 267--269); def of H special cases (lines 271--273); def of q scaling (lines 343--345); Hsummandasymp (lines 1459--1463). Every definition, display and label of those lines is retained unchanged. Omitted from this package and not redistributed: 1--266, 270--270, 274--342, 346--1458, 1464--1519, 1584--1927. Nothing inside a retained excerpt is omitted or altered: every retained body is delivered exactly as the article's lines are written, so no omission receipt of any kind accompanies this package. Citations: the retained excerpts carry no citation command at all, so no bibliography entry of the article is transcribed and no reference is redistributed. Reference adaptation: none was needed and none was made; every reference inside the delivered unit resolves inside it, so no unresolved reference or citation is produced. Printed slips kept literal: no printed word, symbol, hypothesis or number of the counted statement or of its proof was altered. Layout and class: the article's own class is \texttt{amsart} with packages including amssymb, amsthm, amsfonts, amsmath, pmboxdraw, verbatim, graphicx, color, hyperref, cite, ulem, subcaption, bbm, bm; the wrapper is the lane's pinned \texttt{amsart} editorial wrapper at 10pt on A4, which restates the theorem-like environments the retained text needs and adds the article's own macro definitions that the retained text needs (none), with extra wrapper packages (bm, bbm, dsfont, xspace, cancel, mathtools). The delivered numbering is the unit's own. Assets: no retained span contains a figure environment, a graphics-inclusion command, a PDF-page inclusion, a table or a TikZ picture; no image file is copied into this package, the package declares no asset dependency and no image file is redistributed with it. Licence: version arXiv:2507.18042v1 of this article is distributed under the Creative Commons Attribution 4.0 International licence, as recorded in the arXiv licence metadata for this exact version and shown on the abstract page https://arxiv.org/abs/2507.18042v1. A full rights notice is delivered at licenses/arxiv-2507.18042-CC-BY-4.0.txt. Characters: every retained excerpt is pure ASCII, so the delivered engine, XeLaTeX with its default Latin Modern text font, reports no missing character.
    \begin{equation}\label{def of H}
        \mathsf{H}(b,c) :=\sum_{0\leq j_{1}\leq j_{2}\leq\cdots\leq j_{c}\leq b}\prod_{k=1}^{c}\frac{[2j_{k}+k-2]_{q}!!}{[2j_{k}+k-1]_{q}!!}.
    \end{equation} 

\begin{equation}\label{def of H special cases}
\mathsf{H}(b,0)=1, \qquad \mathsf{H}(0,c)=\frac{1}{[c-1]_{q}!!}.
\end{equation}

\begin{equation} \label{def of q scaling}
    q=e^{-\frac{\lambda}{N}}, \qquad \lambda\geq0.
\end{equation} 

    \begin{equation}\label{Hsummandasymp}
        \prod_{r=1}^{2p-2k}\frac{[2j_{r}+r-2]_{q}!!}{[2j_{r}+r-1]_{q}!!}
        =\prod_{r=1}^{2p-2k}\frac{(2j_r +r-2)!!}{(2j_r+r-1)!!}
        \Big(1+\mathcal{D}_{k,1}\frac{\lambda}{N}+\mathcal{D}_{k,2}\frac{\lambda^2}{N^{2}}+O(N^{-3})\Big), 
    \end{equation}

\begin{lem}\label{lem:H asymptotic}
    Let $q$ be scaled according to \eqref{def of q scaling}. 
Let $\mathsf{H}(b,c)$ be given by \eqref{def of H}.  Then as $N\rightarrow\infty$ we have
    \begin{align}
        \begin{split}
        \mathsf{H}(0,2p-2l)
        &=\frac{1}{(2p-2l-1)!!}\bigg(1+\frac{(p-l)(p-l-1)}{2}\frac{\lambda}{N}\\
        &\quad+\frac{(p-l)(p-l-1)(-4+13l+9l^2-13p-18lp+9p^2)}{72}\frac{\lambda^2}{N^2}+O(N^{-3})\bigg)\label{eqn:H(0,2p-2l)},
        \end{split}
        \\
        \mathsf{H}(1,2p-2l-2)
        &=\frac{p-l}{(2p-2l-3)!!}\bigg(1+\frac{(p-l-1)(3p-3l-4)}{6}\frac{\lambda}{N}+O(N^{-2})\bigg), 
    \label{eqn:H(1,2p-2l-2)}
    \\
    \label{eqn:H(2,2p-2l-4)}
        \mathsf{H}(2,2p-2l-4)& =\frac{(p-l)(p-l-1)}{2(2p-2l-5)!!}+O(N^{-1}).
    \end{align} 
\end{lem}

\begin{proof}
    Note that by \eqref{def of H special cases}, we have
    \begin{equation*}
        \mathsf{H}(0,2p-2l)=\frac{1}{[2p-2l-1]_{q}!!}.
    \end{equation*}
    Then \eqref{eqn:H(0,2p-2l)} follows from straightforward computations. 
    
    Next, we show  \eqref{eqn:H(1,2p-2l-2)}. 
    It follows from \eqref{Hsummandasymp} that
    \begin{align}
    \begin{split}
    \label{eqn:asymptotic of H(1,2p-2l-2)}
        \mathsf{H}(1,2p-2l-2)&=\sum_{0\leq j_{1}{\leq}\cdots{\leq}j_{2p-2l-2}\leq1}
        \prod_{r=1}^{2p-2l-2}\frac{(2j_{r}+r-2)!!}{(2j_{r}+r-1)!!}
        \\
        &\qquad \times \bigg(1+\Big(\frac{1}{2}\sum_{r=1}^{2p-2l-2}j_r +\frac{(p-l-2)(p-l-1)}{2}\Big)\frac{\lambda}{N}+O(N^{-2})\bigg).
    \end{split}
    \end{align} 
    Let $m$ be the number of indices $c$ such that $j_c=0$. Then we have
    \begin{align*}
        &\quad \sum_{0\leq j_{1}{\leq} \cdots{\leq}j_{2p-2l-2}\leq1}\prod_{r=1}^{2p-2k-2}\frac{(2j_{r}+r-2)!!}{(2j_{r}+r-1)!!}
        =\sum_{m=0}^{2p-2l-2}
        \Big(\prod_{r=1}^{m}\frac{(r-2)!!}{(r-1)!!}\Big)
        \Big(\prod_{r=m+1}^{2p-2l-2}\frac{r!!}{(r+1)!!}\Big)
        \\
        &=\sum_{m=0}^{2p-2l-2}
        \Big(\prod_{r=1}^{m}\frac{(r-2)!!}{(r-1)!!}\Big)
        \Big(\prod_{r=m+3}^{2p-2l}\frac{(r-2)!!}{(r-1)!!}\Big)= \sum_{m=0}^{2p-2l-2} (m+1)\prod_{r=1}^{2p-2l}\frac{(r-2)!!}{(r-1)!!} =  \sum_{m=0}^{2p-2l-2}  \frac{m+1}{(2p-2l-1)!!}. 
    \end{align*} 
    On the other hand, from the choice of the integer $m$, we also have
    \begin{equation*}
        \sum_{j=1}^{2p-2l-2}j_{r}=2p-2l-2-m. 
    \end{equation*}
    Combining all of the above, we obtain 
    \begin{equation*}
        \mathsf{H}(1,2p-2l-2)
        =\sum_{m=0}^{2p-2l-2}\frac{m+1}{(2p-2l-1)!!}\bigg(1+\Big(\frac{2p-2l-2-m}{2}+\frac{(p-l-2)(p-l-1)}{2}\Big)\frac{\lambda}{N}+O(N^{-2})\bigg),
    \end{equation*}
    which leads to \eqref{eqn:H(1,2p-2l-2)} after simplifications. 
    
   Finally, for \eqref{eqn:H(2,2p-2l-4)}, we only require the leading-order term. Thus, it suffices to establish the identity
    \begin{equation}
        \sum_{0\leq j_1\leq\cdots\leq j_{n}\leq 2}\prod_{r=1}^{n}\frac{(2j_r +r-2)!!}{(2j_r +r-1)!!}=\frac{n^4+10n^3+35n^2+50n+24}{8(n+3)!!}
    \end{equation}
    for $n\geq1$. This summation can be evaluated using a similar strategy to that used in the derivation of \eqref{eqn:H(1,2p-2l-2)}, by decomposing the sum according to the number of indices \( c \) for which \( j_c = 0 \), \( j_c = 1 \), or \( j_c = 2 \). This completes the proof. 
\end{proof}

\end{document}
